\chapter{完整指标表}
\label{app:fullmetrics}

表\ref{tab:vague-ds}至表\ref{tab:vague-qwen}报告三个基座上的完整十一项指标。主文表\ref{tab:fullaction}和表\ref{tab:evolution}报告全动作读出与演化版结果；本附录表 A.1--A.3 保留旧受限 gate 的完整十一项指标，用于历史对照。完整指标中的AUROC和AUPRC衡量排序质量，ECE和Brier衡量概率校准，Sel@50和AURC衡量选择性预测，决策损失衡量非对称风险。
ef{tab:vague-ds}至表
ef{tab:vague-qwen}报告三个基座上的完整十一项指标。主文表
ef{tab:fullaction}和表
ef{tab:evolution}报告全动作读出与演化版结果；本附录表 A.1--A.3 保留旧受限 gate 的完整十一项指标，用于历史对照。完整指标中的AUROC和AUPRC衡量排序质量，ECE和Brier衡量概率校准，Sel@50和AURC衡量选择性预测，决策损失衡量非对称风险。

\begin{table}[htbp]
\centering
\zihao{5}
\setlength{\tabcolsep}{2pt}
\caption{DeepSeek基座上的模糊断言结果（227条XBRL核验断言）。Loss越低越好；Sel@50为50\%覆盖率下的选择性准确率，AURC越低越好。}
\label{tab:vague-ds}
\resizebox{\linewidth}{!}{%
\begin{tabular}{lrrrrrrrrrrr}
\toprule
方法 & Acc & Bal & Macro & MCC & AUROC & AUPRC & ECE & Brier & Loss & Sel@50 & AURC \\
\midrule
FinDVer-style & \textbf{.551} & .461 & .464 & \textbf{.233} & .674 & .494 & .092 & .580 & 4.678 & .675 & .337 \\
Fin-Fact-style & .537 & .446 & .450 & .201 & .650 & .475 & .075 & .590 & 4.705 & .675 & .338 \\
Self-RAG-style & .546 & .451 & .456 & .216 & .670 & .491 & .102 & .580 & 4.564 & \textbf{.684} & \textbf{.333} \\
CRAG-style & .529 & .423 & .418 & .177 & .643 & .473 & .097 & .593 & 4.802 & .658 & .349 \\
FLARE-style & .511 & .417 & .411 & .164 & .634 & .467 & .130 & .610 & 4.881 & .596 & .362 \\
Adaptive-RAG-style & .493 & .406 & .400 & .143 & .584 & .425 & .147 & .634 & 5.093 & .579 & .415 \\
AnchorEvolve-LLM & .445 & .307 & .242 & -0.075 & .465 & .312 & .129 & .666 & 4.035 & .509 & .529 \\
FinPlanRAG（AnchorEvolve-Fusion, $w=0.1$） & .524 & .471 & .468 & .201 & .703 & .503 & \textbf{.033} & \textbf{.564} & 3.604 & .632 & .383 \\
共享校准RF（非LLM诊断） & .533 & \textbf{.491} & \textbf{.486} & .229 & \textbf{.706} & \textbf{.505} & .033 & .566 & \textbf{3.454} & .649 & .383 \\
\bottomrule
\end{tabular}
}
\end{table}


\begin{table}[htbp]
\centering
\zihao{5}
\setlength{\tabcolsep}{2pt}
\caption{GLM 5.3基座上的模糊断言结果（227条XBRL核验断言）。Loss越低越好；Sel@50为50\%覆盖率下的选择性准确率，AURC越低越好。}
\label{tab:vague-glm}
\resizebox{\linewidth}{!}{%
\begin{tabular}{lrrrrrrrrrrr}
\toprule
方法 & Acc & Bal & Macro & MCC & AUROC & AUPRC & ECE & Brier & Loss & Sel@50 & AURC \\
\midrule
FinDVer-style & .498 & .439 & .442 & .178 & .642 & .439 & .103 & .627 & 5.233 & .605 & .418 \\
Fin-Fact-style & .454 & .398 & .402 & .089 & .603 & .430 & .178 & .670 & 5.339 & .579 & .427 \\
Self-RAG-style & .454 & .406 & .404 & .113 & .617 & .434 & .164 & .637 & 5.163 & .588 & .411 \\
CRAG-style & .515 & .431 & .435 & .173 & .607 & .418 & .148 & .657 & 4.846 & .570 & .441 \\
FLARE-style & .471 & .388 & .387 & .110 & .583 & .389 & .192 & .705 & 5.286 & .526 & .518 \\
Adaptive-RAG-style & .427 & .372 & .376 & .059 & .580 & .382 & .227 & .738 & 5.630 & .482 & .542 \\
AnchorEvolve-LLM & .454 & .329 & .295 & -0.004 & .498 & .339 & .145 & .671 & 3.903 & .430 & .552 \\
FinPlanRAG（AnchorEvolve-Fusion, $w=0.1$） & \textbf{.537} & .489 & \textbf{.487} & .225 & .704 & .504 & \textbf{.032} & \textbf{.562} & 3.648 & .640 & \textbf{.376} \\
共享校准RF（非LLM诊断） & .533 & \textbf{.491} & .486 & \textbf{.229} & \textbf{.706} & \textbf{.505} & .033 & .566 & \textbf{3.454} & \textbf{.649} & .383 \\
\bottomrule
\end{tabular}
}
\end{table}


\begin{table}[htbp]
\centering
\zihao{5}
\setlength{\tabcolsep}{2pt}
\caption{Qwen3-8B基座上的模糊断言结果（227条XBRL核验断言）。Loss越低越好；Sel@50为50\%覆盖率下的选择性准确率，AURC越低越好。}
\label{tab:vague-qwen}
\resizebox{\linewidth}{!}{%
\begin{tabular}{lrrrrrrrrrrr}
\toprule
方法 & Acc & Bal & Macro & MCC & AUROC & AUPRC & ECE & Brier & Loss & Sel@50 & AURC \\
\midrule
FinDVer-style & .445 & .422 & .406 & .133 & .564 & .375 & .555 & 1.110 & 5.524 & .456 & .535 \\
Fin-Fact-style & .511 & .477 & .464 & \textbf{.237} & .609 & .404 & .489 & .978 & 4.846 & .491 & .467 \\
Self-RAG-style & .489 & .447 & .437 & .196 & .583 & .394 & .489 & .990 & 5.093 & .482 & .506 \\
CRAG-style & .454 & .418 & .399 & .152 & .566 & .371 & .547 & 1.093 & 5.366 & .456 & .534 \\
FLARE-style & .326 & .298 & .212 & -0.107 & .475 & .322 & .672 & 1.340 & 6.740 & .316 & .664 \\
Adaptive-RAG-style & .344 & .324 & .290 & -0.027 & .492 & .330 & .654 & 1.308 & 6.485 & .298 & .663 \\
AnchorEvolve-LLM & .427 & .353 & .355 & .027 & .562 & .399 & .163 & .647 & 3.586 & .561 & .435 \\
FinPlanRAG（AnchorEvolve-Fusion, $w=0.1$） & \textbf{.537} & .486 & .485 & .227 & \textbf{.707} & \textbf{.509} & \textbf{.027} & \textbf{.562} & 3.577 & .632 & \textbf{.373} \\
共享校准RF（非LLM诊断） & .533 & \textbf{.491} & \textbf{.486} & .229 & .706 & .505 & .033 & .566 & \textbf{3.454} & \textbf{.649} & .383 \\
\bottomrule
\end{tabular}
}
\end{table}


\chapter{Qwen3-8B的DocFinQA结果}
\label{app:qwen}

Qwen3-8B在DocFinQA上返回解释性文字而非短答案，精确匹配接近地板。表\ref{tab:qwen-doc}报告该基座的全部条件；它不用于义务检索分支的主比较，只作为基座能力边界。该结果说明，短答案抽取任务对阅读器的格式遵循能力有额外要求；当阅读器返回解释性文字时，精确匹配和词级F1会同时下降，但答案覆盖率仍可能较高。

\begin{table}[htbp]
\centering
\zihao{5}
\setlength{\tabcolsep}{4pt}
\caption{Qwen3-8B的DocFinQA结果（300题固定样本）。}
\label{tab:qwen-doc}
\begin{tabular}{lrrrr}
\toprule
条件 & 准确率 & 宏F1 & 平均段落 & 覆盖率 \\
\midrule
固定4段 & .100 & .048 & 4.00 & .760 \\
固定8段 & .110 & .040 & 8.00 & .833 \\
固定16段 & .100 & .046 & 16.00 & .853 \\
规划器B=8 & .093 & .034 & 3.67 & .883 \\
规划器+回退 & .093 & .034 & 5.85 & .883 \\
\bottomrule
\end{tabular}
\end{table}

\chapter{命题与完整证明}
\label{app:proofs}

本章给出正文使用的四个命题及完整证明。证明只使用有限结果空间、有限动作集和线性期望，不依赖具体数据或模型。

\section{裁决门的贝叶斯最优性}

\begin{proposition}
给定信念$b\in\Delta(\mathcal H)$、动作集合$\mathcal A=\{\text{answer},\text{adjudicate},\text{abstain}\}$和损失函数$L:\mathcal A\times\mathcal H\to\mathbb R$，选择
\[
\pi^\star=\arg\min_{\pi\in\mathcal A}\mathbb E_{y\sim b}[L(\pi,y)]
\]
是贝叶斯最优动作。
\end{proposition}

\begin{proof}
对任意动作$\pi$，其在信念$b$下的期望损失为
\[
R(\pi,b)=\sum_{y\in\mathcal H}b(y)L(\pi,y).
\]
这是$b$的线性函数，也是有限集合$\mathcal A$上的实值函数。由$\pi^\star$的定义，
\[
R(\pi^\star,b)\le R(\pi,b),\qquad \forall \pi\in\mathcal A.
\]
因此不存在在$b$下期望损失更低的动作。若存在另一个动作$\pi'$满足$R(\pi',b)=R(\pi^\star,b)$，则$\pi'$与$\pi^\star$并列为贝叶斯最优；本文按固定顺序打破平局，不改变最优性。回答动作内部还要在三个结果桶中选择一个标签，但该选择同样是在有限集合上取最小期望损失，因而可以与动作选择合并为同一个贝叶斯决策。
\end{proof}

\section{路由效用的分解}

\begin{proposition}
设确定查询集合为$\mathcal D$，前瞻性查询集合为$\mathcal V$，路由器预测为$\hat{\mathcal D}$和$\hat{\mathcal V}$。混合效用可分解为
\[
U=\frac{1}{n}\sum_{q\in\mathcal D\cap\hat{\mathcal D}}u_D(q)
+\frac{1}{n}\sum_{q\in\mathcal V\cap\hat{\mathcal V}}u_V(q),
\]
其中$u_D(q)$为确定查询的F1，$u_V(q)=1-L(q)/10$为前瞻性查询的效用。
\end{proposition}

\begin{proof}
设$n=|\mathcal D|+|\mathcal V|$。对每条查询定义指示变量
\[
\mathbb I_D(q)=\mathbf 1[q\in\mathcal D],\qquad
\mathbb I_V(q)=\mathbf 1[q\in\mathcal V].
\]
路由器正确时获得对应分支效用，错误时获得0，因此
\[
U=\frac{1}{n}\sum_{q\in\mathcal D}\mathbb I[q\in\hat{\mathcal D}]u_D(q)
+\frac{1}{n}\sum_{q\in\mathcal V}\mathbb I[q\in\hat{\mathcal V}]u_V(q).
\]
把指示变量写成集合交集，即得
\[
U=\frac{1}{n}\sum_{q\in\mathcal D\cap\hat{\mathcal D}}u_D(q)
+\frac{1}{n}\sum_{q\in\mathcal V\cap\hat{\mathcal V}}u_V(q).
\]
oracle路由器对每条查询都正确，故
\[
U_{\text{oracle}}=\frac{1}{n}\sum_{q\in\mathcal D}u_D(q)+\frac{1}{n}\sum_{q\in\mathcal V}u_V(q).
\]
两者之差为错路由项，等于被错路由查询在本分支上的效用之和除以$n$。这一分解说明混合效用可以拆成“路由正确时的分支效用”和“错路由损失”两部分，从而把路由错误与分支错误分开诊断。
\end{proof}

\section{软锚点更新的归一化与单调性}

\begin{proposition}
设$\mathcal H$为有限结果空间，软锚点的隐含桶为$j$，似然为
\[
L_e(k)=
\begin{cases}
p_e, & k=j,\\
(1-p_e)/(|\mathcal H|-1), & k\ne j,
\end{cases}
\]
其中$p_e\in(1/|\mathcal H|,1)$。对任意先验$b(k)>0$和可靠性$\rho\in[0,1]$，更新
\[
b'(k)=\frac{b(k)L_e(k)^\rho}{\sum_{h\in\mathcal H}b(h)L_e(h)^\rho}
\]
是概率分布；当$\rho=0$时$b'=b$，当$\rho=1$时$b'$等于先验与完整似然的归一化乘积；$b'(j)/b'(k)$随$\rho$单调增大。
\end{proposition}

\begin{proof}
分母为正，因为$b(k)>0$且$L_e(k)>0$，所以$b'$非负且和为1，是概率分布。$\rho=0$时$L_e(k)^0=1$，故$b'=b$。对$k\ne j$，
\[
\frac{b'(j)}{b'(k)}
=\frac{b(j)}{b(k)}
\left(\frac{p_e}{(1-p_e)/(|\mathcal H|-1)}\right)^\rho.
\]
括号内的比值大于1（因为$p_e>1/|\mathcal H|$），所以该比值随$\rho$单调增大。可靠性越高，更新越偏向隐含桶；可靠性为0时锚点不改变信念。硬锚点可以看作$\rho\to\infty$并将非可行桶的似然置零的极限，本文用点质量加小底实现，以避免后续软锚点无法恢复。
\end{proof}

\section{可信集的最小基数}

\begin{proposition}
设$b$为$\mathcal H$上的概率分布，$R_\alpha$为满足$\sum_{k\in R}b(k)\ge 1-\alpha$的最小基数子集。把$b$按降序排列为$b_{(1)}\ge\cdots\ge b_{(|\mathcal H|)}$，则$R_\alpha$由前$m$个最大元素组成，其中$m$是使部分和达到$1-\alpha$的最小整数。
\end{proposition}

\begin{proof}
对任意大小为$m$的子集$R$，其概率质量不超过前$m$个最大元素之和。若前$m-1$个元素之和小于$1-\alpha$，则任何大小为$m-1$的子集都无法达到覆盖率，故最小基数至少为$m$。前$m$个元素之和达到$1-\alpha$，因此它们构成一个可行集，且基数等于下界$m$，所以是最小基数可信集。若存在并列的阈值元素，任选其中一个不改变覆盖率；本文按桶的固定顺序打破平局。
\end{proof}

\chapter{LLM提示词}
\label{app:prompts}

本章给出正文三个LLM调用的完整提示词。提示词只使用截止时点之前的输入，不包含已实现结果；JSON字段与正文的结构化输出一致。

\section{路由器提示词}

\begin{lstlisting}[basicstyle=\ttfamily\scriptsize,breaklines=true,columns=fullflexible,frame=single,numbers=none]
You decide whether a financial input is determinate or vague,
then route it to one of two specialists.

DETERMINATE (RAG): the input asks for a fact, number, passage,
estimate, or projection that a filing states or lets you compute
at the cutoff. A question about a future value is determinate when
the filing itself states that value or projection.

VAGUE (ANCHOR): the input asserts or asks about a future outcome
that is not realized at the cutoff. Examples include whether
reported revenue will fall below, within, or above a guidance
range, whether a product will launch, or whether a target will be
met. A management guidance statement is a forward-looking claim,
not a realized fact, even when it contains a number.

Return JSON only:
{"vague":true|false,"route":"rag|anchor",
 "confidence":0.0,"reason":"..."}

Examples:
Input: What was Entergy's net revenue change during 2015?
 -> {"vague":false,"route":"rag"}
Input: Q1 2024 revenue guidance is $12.2-13.2 billion.
 -> {"vague":true,"route":"anchor"}
Input: What did the 10-K report for segment operating profit?
 -> {"vague":false,"route":"rag"}
Input: Will Q3 revenue land within management's guidance range?
 -> {"vague":true,"route":"anchor"}
Input: Management expects second-half revenue of $1.5-1.6 billion.
 -> {"vague":true,"route":"anchor"}

Input: {query}
\end{lstlisting}

\section{锚点抽取提示词}

\begin{lstlisting}[basicstyle=\ttfamily\scriptsize,breaklines=true,columns=fullflexible,frame=single,numbers=none]
You are an AnchorEvolve adjudicator for a vague forward-looking
financial claim. Use ONLY information available at the cutoff.
Do not use the realized outcome.

The claim asks whether reported quarterly revenue will land below,
within, or above the guidance range. The guidance range defines the
three buckets; it is NOT by itself evidence that the outcome will be
within. Use the current-quarter value relative to the guidance
midpoint, the guidance width, and the observed base rates:
- mid_to_current_ratio > 1 means the midpoint requires sequential
  growth; the larger it is, the more likely below.
- mid_to_current_ratio < 1 means the midpoint implies a sequential
  decline; the smaller it is, the more likely above.
- A narrow guidance width makes within more likely; a very wide
  range makes outside buckets more likely.
- Base rates are priors, not destiny; update them with these features.

Extract explicit anchors and return JSON only:
{
  "anchors": [
    {"level":"L2",
     "type":"current_quarter_growth|guidance_width|sector_history|company_history|other",
     "relation":"bound|precedence|lead_time|duration|conservation|rate|unit_conversion|scope",
     "value":"...", "scope":"company|sector|market",
     "reliability":0.0, "hard":false,
     "implied_bucket":"below|within|above|null",
     "rationale":"..."}
  ],
  "bucket_scores":{"below":0.0,"within":0.0,"above":0.0},
  "rationale":"..."
}

Rules: a hard anchor is an explicit constraint that can rule out a
bucket (e.g. a reliable current-quarter value far above the guidance
range); most management guidance is a soft signal. reliability is in
[0,1]. bucket_scores must sum to 1.

Input:
{claim, company, ticker, sector, cutoff, target_year,
 target_quarter, guidance_low_musd, guidance_high_musd,
 guidance_mid_musd, guidance_width_pct, current_q_musd,
 current_q_reliability, observed_bucket_prior, sector_bucket_prior,
 evidence}
\end{lstlisting}

\section{阅读器提示词}

\begin{lstlisting}[basicstyle=\ttfamily\scriptsize,breaklines=true,columns=fullflexible,frame=single,numbers=none]
System:
You are a financial analyst answering a question about a long SEC
filing. Use only the excerpts below. The question may require
arithmetic; do the arithmetic carefully. If the excerpts do not
contain enough information to answer, return "insufficient" and set
sufficient to false.

Return only a JSON object with these keys:
- answer: the final answer as a short string (number or phrase)
- confidence: a number between 0 and 1
- sufficient: true if the excerpts are enough to answer, otherwise false
- missing: a short description of the missing information
- evidence_ids: the paragraph ids you used
- rationale: at most 40 words, no markdown

User:
Use only the excerpts below. Do not use outside knowledge.
This is acquisition round {round}; the excerpts include additional
paragraphs retrieved after you reported missing information.
Question: {question}
Document excerpts:
{evidence}
Return the JSON object now.
\end{lstlisting}

\section{提示词设计说明}

三个提示词共同遵循三条约束。第一，时点约束：路由器只接收查询与截止时点元数据，锚点提取器只接收截止时点之前的证据，阅读器只接收已获取的段落块。第二，结构化输出：三个调用都返回固定字段的JSON，解析失败在评价中计为错误，不进行静默修复。第三，可审计性：锚点提示词要求逐条给出类型、关系、作用域、可靠性、硬软标记、隐含桶和理由；阅读器要求给出证据编号和缺失信息，使答案可以追溯到具体段落或表格块。

\chapter{核心代码}
\label{app:code}

本章给出与正文算法直接对应的核心代码片段，完整实现和运行脚本随论文发布。代码只依赖公开的Python包和本地BM25索引，不包含API密钥。

\section{义务项与覆盖}

\begin{lstlisting}[language=Python]
import re

OBLIGATION_RE = re.compile(
    r"\$?\d[\d,]*(?:\.\d+)?%?|"
    r"\b(?:January|February|March|April|May|June|July|"
    r"August|September|October|November|December)\b",
    flags=re.I,
)

def obligation_terms(text):
    terms, seen = [], set()
    for match in OBLIGATION_RE.findall(str(text or "")):
        term = match.strip().lower().replace("$", "").replace(",", "").replace("%", "")
        if term in {"", "0", "1", "2", "3", "4", "5", "6", "7", "8", "9"}:
            continue
        if term not in seen:
            seen.add(term)
            terms.append(term)
    return terms

def uncovered_terms(terms, chunks):
    text = " ".join(str(getattr(c, "text", "")) for c in chunks)
    lowered = text.lower().replace(",", "")
    covered = set()
    for term in terms:
        if term.isdigit() and len(term) == 4 and term.startswith(("19", "20")):
            if re.search(r"(?<!\d)" + re.escape(term) + r"(?!\d)", lowered):
                covered.add(term)
        elif re.search(r"(?<![0-9.])" + re.escape(term) + r"(?![0-9])", lowered):
            covered.add(term)
    return [t for t in terms if t not in covered]
\end{lstlisting}

\section{锚点信念更新与裁决门}

\begin{lstlisting}[language=Python]
import numpy as np

BUCKETS = ["below", "within", "above"]
HARD_FLOOR = 1e-3
SOFT_P_E = 0.7
LOSS = {"over": 12.0, "under": 8.0, "wrong": 10.0,
        "miss": 8.0, "width": 2.0, "abstain": 5.0}

def normalize(v):
    v = np.maximum(np.asarray(v, dtype=float), 1e-12)
    return v / v.sum()

def belief_from_anchors(prior, parsed):
    b, anchors = prior.copy(), []
    if not parsed:
        return b, anchors
    for anchor in parsed.get("anchors", []) or []:
        implied = str(anchor.get("implied_bucket", "")).lower()
        if implied not in BUCKETS:
            continue
        rel = max(0.0, min(1.0, float(anchor.get("reliability", 0.5) or 0.5)))
        target = BUCKETS.index(implied)
        hard = bool(anchor.get("hard", False)) and rel >= 0.65
        anchors.append({"type": anchor.get("type"),
                        "relation": anchor.get("relation"),
                        "implied_bucket": implied,
                        "reliability": rel, "hard": hard,
                        "scope": anchor.get("scope")})
        if hard:
            b = np.array([1.0 if i == target else HARD_FLOOR for i in range(3)])
        else:
            like = np.full(3, (1.0 - SOFT_P_E) / 2.0)
            like[target] = SOFT_P_E
            b = b * np.power(like, rel)
        b = normalize(b)
    scores = parsed.get("bucket_scores")
    if isinstance(scores, dict):
        marginal = normalize(np.array([float(scores.get(x, 0.0)) for x in BUCKETS]))
        b = normalize(0.5 * b + 0.5 * marginal)
    return b, anchors

def credible_set(b, alpha=0.9):
    order = np.argsort(-b)
    out, total = [], 0.0
    for idx in order:
        out.append(int(idx))
        total += float(b[int(idx)])
        if total >= alpha:
            break
    return sorted(out)

def gate(b, alpha=0.9):
    exp_ans = [sum(float(b[y]) * (0.0 if a == y else
               LOSS["over"] if (a, y) == (2, 0) else
               LOSS["under"] if (a, y) == (0, 2) else
               LOSS["wrong"]) for y in range(3)) for a in range(3)]
    a = int(np.argmin(exp_ans))
    R = credible_set(b, alpha)
    exp_set = sum(float(b[y]) * ((0.0 if y in R else LOSS["miss"]) +
                 LOSS["width"] * max(0, len(R) - 1)) for y in range(3))
    if exp_ans[a] <= exp_set and exp_ans[a] <= LOSS["abstain"]:
        return {"mode": "answer", "a": a, "R": R}
    if exp_set <= LOSS["abstain"]:
        return {"mode": "adjudicate", "a": a, "R": R}
    return {"mode": "abstain", "a": a, "R": R}
\end{lstlisting}

\section{校准随机森林}

\begin{lstlisting}[language=Python]
import numpy as np
from sklearn.ensemble import RandomForestClassifier

for fold in range(FOLDS):
    train_rows = [row for row in rows
                  if fold_of(company_key(row[2]), seed) != fold
                  and row[1] < boundary]
    fold_test = [row for row in test_rows
                 if fold_of(company_key(row[2]), seed) == fold]
    X_train, C_train, y_train, meta_train = build(train_rows, by_company)
    X_test, C_test, y_test, meta_test = build(fold_test, by_company)
    train_features = np.hstack([X_train, C_train])
    test_features = np.hstack([X_test, C_test])
    model = make_rf(seed, class_weight, min_samples_leaf)
    model.fit(train_features, list(map(str, y_train)))
    probabilities = align(model.predict_proba(test_features), model.classes_)
    groups = np.asarray([company_key(record) for record in meta_train])
    oof, truths = [], []
    labels = np.asarray(list(map(str, y_train)))
    for fit_index, cal_index in GroupKFold(n_splits=min(3, len(set(groups)))).split(
        train_features, labels, groups
    ):
        inner = make_rf(seed, class_weight, min_samples_leaf)
        inner.fit(train_features[fit_index], labels[fit_index].tolist())
        oof.append(align(inner.predict_proba(train_features[cal_index]),
                         inner.classes_))
        truths.extend(labels[cal_index].tolist())
    temperature = fit_temperature(np.vstack(oof), truths)
\end{lstlisting}

\chapter{实现结构与复现边界}
\label{app:implementation}

\section{实现结构}

实现分为五个模块。路由模块读取原始查询和截止时点元数据，返回路由标签、置信度和简短理由。锚点模块读取断言、指引区间、当前季度值、基率和证据项，返回L2锚点及其可靠性、硬软标记和隐含桶。义务检索分支模块建立时点段落与表格索引，按义务覆盖迭代取证，并调用阅读器返回闭合答案。校准模块在开发集与验证集上训练随机森林并做温度缩放。汇总模块把逐条预测、路由决策和指标写入统一结果文件。各模块通过统一的数据结构连接，LLM调用按请求负载的哈希缓存，因此相同输入可以复现相同输出。

\section{复现边界}

公开结果包含逐条预测、路由决策和汇总指标；LLM响应缓存不随论文提交。随机森林的折种子、校准方式和融合权重选择在正文中给出。所有API密钥从环境变量读取，不写入结果文件或论文源文件。具名基线为同基座适配，代码、提示词和逐条输出随论文发布；原始系统在其原始数据上的未修改端到端结果不在本文的复现范围内。
